\chapter{Preface and reading guide}

This manuscript studies one question: how can analytic control of quadratic
dynamics produce connected neighbourhoods in the Mandelbrot parameter
plane?  The mathematical account and the Lean development are presented
together, but their completion levels are recorded separately.  Published
theorems remain mathematical theorems before their proofs are formalized;
conversely, a checked Lean implication represents a published theorem only
after its definitions and hypotheses have been identified with the
objects in the source.

The exposition follows the progression of Lyubich's conformal-geometry
manuscripts~\cite{LyubichBook,LyubichNotes}: elementary dynamics, topology,
external coordinates, puzzles, and renormalization.  The categorical
approaches are discussed after those foundations, so that their geometric
inputs can be compared with the literature.  Basic real and complex
analysis are prerequisites; the dynamical objects, neighbourhood
conditions, and renormalization terminology used in the argument are
introduced in the text.

\section*{Four routes through the book}
\addcontentsline{toc}{section}{Four routes through the book}

Part~I defines the quadratic family and the local-connectivity target,
proves the escape criterion, and develops the compactness and component
arguments used by every later route.  Part~II introduces external
coordinates and the phase--parameter problem before discussing
renormalization and its discrete bookkeeping.  Part~III compares the
published parameter classes and separates direct MLC conclusions from
supporting geometric and rigidity results.  Part~IV audits the current
formal models and then explains the conditional categorical and
certificate-based roots.

The appendices are reference material: a dependency map, a declaration
inventory, and the source policy.  File numbers in
\lean{docs/review/} record the history of the manuscript; chapter numbers
and cross-references follow the mathematical order in the master file.
Thus no mathematical dependency should be inferred from a filename.

\section*{Reading a theorem together with its code}
\addcontentsline{toc}{section}{Reading a theorem together with its code}

The book uses the following distinctions throughout.

\begin{center}\small
\begin{tabularx}{\textwidth}{@{}P{0.25\textwidth}Y@{}}
\textbf{Status} & \textbf{Meaning}\\
\hline
Kernel-checked &
The exact Lean proposition follows from its stated hypotheses and audited
dependencies.  The meaning of every definition still matters.\\
Literature theorem &
A cited paper proves the mathematical result.  The corresponding Lean
proof may remain incomplete.\\
Conditional interface &
The code checks the implication from supplied data.  Constructing those
data is a separate task.\\
Formalization obligation &
A construction or proof still has to be implemented.  Within a published
parameter class this may already be known mathematics.\\
Open mathematical problem &
The cited literature supplies no proof in the stated generality.\\
Refuted model input &
A counterexample or a general obstruction prevents the proposed data from
existing with their present definitions.\\
\end{tabularx}
\end{center}

Each substantial bridge is stated mathematically before its Lean
correspondence is given.  A linked declaration name opens the containing
source file on the working branch.  The file and the declaration name
together identify the formal statement; the pinned dependency revisions
are recorded by \lean{lake-manifest.json}.  Long identifiers are allowed to
break across lines rather than widening the text block.

\begin{conditionalstep}
The present audit finds an important mismatch in the inherited puzzle
model.  Its declaration \lean{Quadratic.modulus} is the integral of a
Gaussian weight, whereas the literature uses conformal annulus modulus.
The former is denoted by $\weightedarea$ in this book.  Its values on
nested differences always form a summable series.
Chapter~\ref{chap:frozen-audit} proves and formalizes this obstruction.
Consequently, the legacy theorem named
\lean{yoccoz_puzzle_shrinkage} is a checked implication with an impossible
premise, rather than a completed formalization of Yoccoz's theorem.
\end{conditionalstep}

\section*{Notation}
\addcontentsline{toc}{section}{Notation}

\begin{center}\small
\begin{tabularx}{\textwidth}{@{}P{0.28\textwidth}Y@{}}
\textbf{Symbol} & \textbf{Meaning}\\
\hline
$f_c,\ p_n(c)$ & $z\mapsto z^2+c$, and its critical orbit
$f_c^{\circ n}(0)$.\\
$K_c,\ J_c,\ \M$ & Filled Julia set, its boundary, and the bounded
critical-orbit locus.\\
$B(x,r),\ \overline B(x,r)$ & Open and closed Euclidean balls.\\
$\overline A,\ \partial A,\ \operatorname{int}A$ &
Closure, boundary, and interior in the stated ambient space.\\
$\mathcal N_X(x)$ & Neighbourhood filter in the space $X$.\\
$\operatorname{Comp}_x(A)$ & Connected component of $A$ containing $x$,
when $x\in A$.\\
$O_N$ & Compact finite-orbit outer approximation of $\M$.\\
$G_c,\ \phi_c,\ \Phi$ & Dynamical Green function, dynamical Böttcher
coordinate, and parameter external uniformization.\\
$P_n(c,z),\ \mathcal P_n(c)$ & Graph-cut dynamical and parameter puzzle
pieces; these symbols refer to the literature construction.\\
$\modulus(A)$ & Conformal modulus of a ring domain.\\
$\weightedarea(A)$ & Gaussian weighted area in the legacy Lean model.\\
$U\Subset V$ & $\overline U$ is compact and contained in $V$.\\
$d_n,\ T_n$ & Additive bookkeeping depth and multiplicative original
return period, respectively.\\
\end{tabularx}
\end{center}

\section*{Manuscripts, papers, and version boundaries}
\addcontentsline{toc}{section}{Manuscripts, papers, and version boundaries}

The author's publication page lists the online book as volumes I--II.
The repository also contains a distinct 2013 ICTP manuscript under the
same main title.  The article \emph{Dynamics of quadratic polynomials,
III: parapuzzle and SBR measures} is a research paper, rather than a third
volume of that book.  The ``IV'' in Kahn--Lyubich's \emph{Elephant Eyes}
belongs to the separate \emph{A priori bounds} series.  No further book
volume is inferred from either numbering scheme.

The literature boundary is reviewed as of 30 October 2026.
Dudko's survey~\cite{DudkoMLC} provides the overall renormalization
roadmap; the 2026 papers are cited individually where their hypotheses and
conclusions enter.  The target remains a full book with full formal
coverage.  This working manuscript records the material currently
developed and the remaining obligations, without treating its length,
source inventory, or theorem names as evidence of completion.
